\section{Traditional social savings and the efficient benchmark}
\label{ch:hulten}

\subsection{The accounting statistic}

Social savings holds the observed shipment pattern fixed and asks how much
less it would cost to move the same bundle after an infrastructure
improvement. For a marginal proportional reduction in one edge--mode cost,
homothetic demand and iceberg costs give the welfare elasticity
\[
 H_{klm}=\Xi_{kl,m}.
\]
The expression is attractive because it needs only the baseline value flow.
It does not require solving a counterfactual spatial equilibrium.

This statistic is often described as a partial-equilibrium approximation.
That description is incomplete. Under the maintained efficient benchmark it
is the exact first-order general-equilibrium effect even though trade, routes,
modes, prices, and locations adjust.

\begin{proposition}[Efficient social savings]
\label{prop:hulten}
Suppose \(\alpha=\beta=0\) and all congestion elasticities are zero. At an
interior efficient equilibrium,
\[
 -\frac{d\log W}{d\log\bar\kappa_{kl,m}}=\Xi_{kl,m}.
\]
\end{proposition}

\subsection{Envelope logic}

The result is an application of the envelope theorem and the network version
of Hulten's theorem \citep{hulten1978growth,baqaee2019macroeconomic}. At an
efficient allocation, the planner is locally indifferent to reallocations
that preserve feasibility. A lower transport cost directly saves resources.
The induced changes in routes, modes, trade, and location are responses along
the feasible set. Their first-order welfare effects vanish at the optimum.

The statement is stronger than saying that adjustment is quantitatively
small. Adjustment can be large. What vanishes is its first-order contribution
to welfare at an efficient baseline. This distinction explains why adding a
rich spatial equilibrium to the model does not, by itself, invalidate social
savings.

\subsection{When the result fails}

Congestion and spatial externalities make the decentralized allocation
inefficient. Then a transport improvement changes the exposure of agents to
pre-existing wedges. Route diversion can relieve or worsen congestion.
Relocation can raise productivity through agglomeration or affect amenities
through dispersion. These responses need not vanish at first order.

The result also depends on the welfare and mobility closure. The paper
formally proves it under perfect labor mobility and the stated aggregate
welfare criterion. Similar envelope logic can apply under other mobility
settings only when the allocation is efficient in the enlarged problem and
welfare is aggregated consistently. Fixed factors, taxes, heterogeneous
households, or endogenous labor supply can add incidence terms.

\subsection{A diagnostic benchmark}

The package retains the Hulten statistic even when the full model is
distorted. It serves three purposes:
\begin{enumerate}
  \item a transparent traffic-based benchmark;
  \item an exact unit test in the efficient limit;
  \item a basis for decomposing why the extended result differs.
\end{enumerate}

If the efficient fixture does not collapse to traffic, the calculation has a
sign, normalization, or indexing error. If the distorted model remains close
to Hulten, that is an empirical result rather than a maintained theorem.

\implementationbox{The output column \code{hulten} is the policy-mode traffic
share. The extended realized and primitive elasticities are reported
separately. Physical-link results sum opposite directions before computing
link-level comparisons.}
